\documentclass[12pt]{article}

\usepackage[T1]{fontenc}
\usepackage[utf8]{inputenc}
\usepackage[margin=1in]{geometry}
\usepackage{amsmath}
\usepackage{booktabs}
\usepackage{longtable}
\usepackage{array}
\usepackage{hyperref}

\title{Supplementary Material for \\ \emph{Credible on Paper? Green Banking Policy, Access to Finance, and Institutional Capacity in the Greening of Firms Worldwide}}
\author{}
\date{}

\begin{document}
\maketitle

This file reports supplementary tables referenced in the main text: a small-sample supplementary
check on a narrower outcome (Table~S1), additional robustness checks and their multiple-testing
correction (Tables~S2--S3), the full per-economy SBFN and regulatory-quality coding underlying the
primary sample (Table~S4), Bayesian convergence diagnostics (Section~S5), and the global sample's
SBFN membership by region (Table~S6). Section and table numbers here (S1, S2, \ldots) are internal
to this supplementary file and distinct from the main text's own numbering, which this file's
prose refers to descriptively (e.g., ``the main text's baseline logit table'') rather than by
number, since table numbers in the main text are subject to change on revision.

\section{Small-sample supplementary check: waste-minimization adoption}
\label{sec:s-extension}

Before the global-sample replication that is the paper's second major test, we report a smaller
supplementary check using waste-minimization adoption---an outcome not captured in the standardized
cross-country database underlying the global sample, and so a genuinely independent (if
small-sample) data point---available in four economies whose Green Economy modules asked this item
under directly comparable wording: India 2022, Indonesia 2023, Peru 2023, and Timor-Leste 2021.
Three of these four are themselves SBFN members at various progression stages, so this sample
cannot re-test the SBFN-adoption contrast; what it can test is whether the credit-to-green channel
documented in the main text's baseline, hierarchical, and causal-forest results is a peculiarity of
the ECA-MENA sample or generalizes elsewhere. It generalizes: pooled across the four economies,
having an overdraft facility is associated with significantly higher odds of waste-minimization
adoption ($b=0.422$, $p<0.001$, $N=9{,}038$, 4 countries). The direction of the core finance-green
relationship replicates on a different continent, a different half-decade, and a narrower outcome
definition than the one anchoring the main text's primary-sample results---a preview, on four
economies, of the much larger replication in the global sample.

\begin{table}[htbp]
\centering
\caption{Small-sample supplementary check: descriptive rates.}
\label{tab:s-extension}
\begin{footnotesize}
\begin{tabular}{lcccc}
\toprule
Economy & $N$ & Waste-min.\ rate & Overdraft rate & SBFN member \\
\midrule
India 2022 & 9{,}376 & 0.515 & 0.621 & 1 \\
Indonesia 2023 & 2{,}955 & 0.336 & 0.144 & 1 \\
Peru 2023 & 987 & 0.806 & 0.549 & 1 \\
Timor-Leste 2021 & 238 & 0.061 & 0.089 & 0 \\
\bottomrule
\end{tabular}
\end{footnotesize}
\end{table}

\section{Additional robustness checks}
\label{sec:s-robustness}

Table~\ref{tab:s-robustness} subjects the primary-sample finding to four further checks, all
retaining sector fixed effects, firm controls, and country-clustered standard errors (no country
fixed effects, for the reason given in the main text). First, replacing the binary outcome with the
continuous seven-item green-adoption index and estimating by OLS reproduces the same pattern to the
decimal: credit access positive and significant ($b=0.053$, $p<0.001$), SBFN main effect negative
and significant ($b=-0.092$, $p<0.001$), interaction indistinguishable from zero ($b=0.006$,
$p=0.739$).

Second, replacing bank credit access with two alternative finance measures produces one genuine
point of nuance worth reporting rather than smoothing over. Using overdraft-facility access instead
of a credit line/loan, the main overdraft effect remains positive and significant ($b=0.274$,
$p=0.011$), the SBFN main effect remains negative and significant ($b=-0.905$, $p<0.001$)---but the
overdraft $\times$ SBFN interaction is now negative and significant ($b=-0.346$, $p=0.031$), the one
interaction term in the entire paper that clears a conventional significance threshold, and it runs
opposite to the direction H2 predicts. Using the reverse-coded perceived finance-obstacle measure in
place of realized access, by contrast, both the main effect and its SBFN interaction are precisely
zero ($b=-0.010$, $p=0.827$; interaction $b=-0.012$, $p=0.849$)---consistent with the descriptive
correlation noted in the main text. We read the overdraft result as suggestive rather than
conclusive given it is one significant coefficient among more than a dozen interaction tests across
the paper (see Table~\ref{tab:s-fdr} below), but it is at least directionally consistent with a
substantive story worth flagging for future work: if SBFN-guided lending is channelled preferentially
through the term loans and credit lines that finance long-horizon capital expenditure, rather than
through short-term overdraft facilities, we might expect exactly this pattern---a null-to-positive
interaction on credit lines and a negative one on overdrafts---though our data cannot directly test
whether SBFN-guided credit is disproportionately term-structured.

Third, splitting the sample into manufacturing ($N=13{,}340$) and services ($N=9{,}538$) subsamples
shows the credit-access main effect, the SBFN main effect, and the null credit $\times$ SBFN
interaction all replicate within each sector separately (manufacturing: credit $b=0.480$, $p<0.001$,
interaction $b=-0.040$, $p=0.826$; services: credit $b=0.448$, $p<0.001$, interaction $b=0.195$,
$p=0.426$)---the paper's central null is not an artifact of pooling two structurally different
production technologies.

\begin{table}[htbp]
\centering
\caption{Additional robustness checks.}
\label{tab:s-robustness}
\begin{footnotesize}
\begin{tabular}{lccc}
\toprule
Specification & Finance coef.\ (p) & SBFN coef.\ (p) & Interaction coef.\ (p) \\
\midrule
(a) Continuous index, OLS & 0.053 (0.000) & -0.092 (0.000) & 0.006 (0.739) \\
(b1) Overdraft & 0.274 (0.011) & -0.905 (0.000) & -0.346 (0.031) \\
(b2) Finance obstacle (reverse-coded) & -0.010 (0.827) & -1.011 (0.000) & -0.012 (0.849) \\
(c1) Manufacturing subsample & 0.480 (0.000) & -1.047 (0.000) & -0.040 (0.826) \\
(c2) Services subsample & 0.448 (0.000) & -0.914 (0.001) & 0.195 (0.426) \\
\bottomrule
\end{tabular}
\end{footnotesize}
\end{table}

Because the overdraft interaction (the one coefficient in the paper that clears a conventional
significance threshold) is one of many credit $\times$ institutional-moderator tests reported across
the frequentist specifications in this paper, we additionally apply a Benjamini--Hochberg
false-discovery-rate correction across the full family of 13 such coefficients reported under a
classical (non-Bayesian, non-forest) estimator in the main text's baseline, robustness, and
global-sample regression tables (Table~\ref{tab:s-fdr} below; full calculation in the replication
package). The overdraft interaction's raw $p=0.031$ becomes $p=0.403$ after correction, and no
coefficient in the family survives at a false discovery rate of 0.10 or even 0.40. We read this as
reinforcing, with a formal multiple-testing accounting rather than only a qualitative caveat, the
suggestive-not-conclusive framing the main text's Discussion already gives this result: an isolated
nominally-significant interaction among more than a dozen tested is exactly what a pure
false-discovery process would be expected to produce roughly one time in twenty, and this correction
does not let us reject that account.

\begin{table}[htbp]
\centering
\caption{Multiple-testing correction: Benjamini--Hochberg false discovery rate applied to every
credit $\times$ institutional-moderator interaction coefficient reported under a frequentist
estimator (main text's baseline, robustness, and global-sample regression tables above). Entries
marked $\dagger$ are approximate $p$-values recovered from the reported coefficient and standard
error via a normal approximation ($z=\text{coef}/\text{se}$), since the main text reports only the
coefficient and standard error for these terms.}
\label{tab:s-fdr}
\begin{footnotesize}
\begin{tabular}{llccc}
\toprule
Source & Term & Raw $p$ & BH-adjusted $p$ \\
\midrule
Table~\ref{tab:s-robustness} (b1) Overdraft & Overdraft $\times$ SBFN & 0.031 & 0.403 \\
Main text, global sample, M2 (country FE) & Credit $\times$ SBFN & 0.205 & 0.799 \\
Main text, global sample, M1 (no country FE) & Credit $\times$ Reg.\ quality$^\dagger$ & 0.301 & 0.799 \\
Main text, baseline M3$^\dagger$ & Credit $\times$ Reg.\ quality & 0.310 & 0.799 \\
Main text, global sample, M2 (country FE) & Credit $\times$ Reg.\ quality & 0.346 & 0.799 \\
Table~\ref{tab:s-robustness} (c2) Services & Credit $\times$ SBFN & 0.426 & 0.799 \\
Main text, global sample, M1 (no country FE) & Credit $\times$ SBFN & 0.430 & 0.799 \\
Main text, baseline M3$^\dagger$ & Credit $\times$ SBFN $\times$ Reg.\ quality & 0.702 & 0.892 \\
Table~\ref{tab:s-robustness} (a) Continuous index, OLS & Credit $\times$ SBFN & 0.739 & 0.892 \\
Main text, baseline M2 & Credit $\times$ SBFN & 0.773 & 0.892 \\
Table~\ref{tab:s-robustness} (c1) Manufacturing & Credit $\times$ SBFN & 0.826 & 0.892 \\
Table~\ref{tab:s-robustness} (b2) Finance obstacle (reverse-coded) & Obstacle $\times$ SBFN & 0.849 & 0.892 \\
Main text, baseline M3$^\dagger$ & Credit $\times$ SBFN & 0.892 & 0.892 \\
\bottomrule
\end{tabular}
\end{footnotesize}
\end{table}

\section{SBFN policy status and regulatory quality, full sample}
\label{app:sbfn}

Table~\ref{tab:appendixa} reports, for all 47 economy-years in the paper's primary and small-sample
supplementary check, the SBFN membership status coded for the analysis, the join year where
applicable, the SBFN Multi-Tiered Progress Framework stage in force nearest the survey year, and the
exact year of the Worldwide Governance Indicators observation used. Full per-country source
citations (the specific SBFN Global Progress Report, country addendum, or accession announcement
consulted for every determination) are available in the replication package. The analogous
162-economy table for the global sample is summarized by region in Table~\ref{tab:globalsbfnregion}
below and provided in full in the replication package.

{\footnotesize
\begin{longtable}{p{2.3cm}ccp{4.0cm}c}
\caption{SBFN policy status and regulatory quality, full 47-economy sample.} \label{tab:appendixa} \\
\toprule
Economy & SBFN member & Join year & Stage at survey & WGI year \\
\midrule
\endfirsthead
\toprule
Economy & SBFN member & Join year & Stage at survey & WGI year \\
\midrule
\endhead
Albania 2019 & 0 & --- & None & 2019 \\
Armenia 2020 & 0 & --- & None & 2020 \\
Azerbaijan 2019 & 0 & --- & None & 2019 \\
Bangladesh 2022 & 1 & 2012 & Advancing (ESG Integration) & 2022 \\
Belarus 2018 & 0 & --- & None & 2018 \\
Bosnia and Herzegovina 2019 & 0 & --- & None & 2019 \\
Bulgaria 2019 & 0 & --- & None & 2019 \\
Croatia 2019 & 0 & --- & None & 2019 \\
Cyprus 2019 & 0 & --- & None & 2019 \\
Czechia 2019 & 0 & --- & None & 2019 \\
Egypt 2020 & 1 & 2016 & Formulating (Preparation) & 2020 \\
Estonia 2019 & 0 & --- & None & 2019 \\
Georgia 2019 & 1 & 2017 & Developing (Implementation) & 2019 \\
Greece 2018 & 0 & --- & None & 2018 \\
Hungary 2019 & 0 & --- & None & 2019 \\
India 2022 & 1 & 2016 & Developing (Commitment) & 2022 \\
Indonesia 2023 & 1 & 2012 & Consolidating (Maturing) & 2023 \\
Italy 2019 & 0 & --- & None & 2019 \\
Jordan 2019 & 1 & 2016 & Commitment (Preparation) & 2019 \\
Kazakhstan 2019 & 0 & --- & None & 2019 \\
Kosovo 2019 & 0 & --- & None & 2019 \\
Kyrgyz Republic 2019 & 1 & 2018 & Commitment (Preparation) & 2019 \\
Latvia 2019 & 0 & --- & None & 2019 \\
Lebanon 2019 & 0 & --- & None & 2019 \\
Lithuania 2019 & 0 & --- & None & 2019 \\
Malta 2019 & 0 & --- & None & 2019 \\
Moldova 2019 & 0 & --- & None & 2019 \\
Mongolia 2019 & 1 & 2012 & Advancing (Implementation) & 2019 \\
Montenegro 2019 & 0 & --- & None & 2019 \\
Morocco 2019 & 1 & 2014 & Advancing (Implementation) & 2019 \\
North Macedonia 2019 & 0 & --- & None & 2019 \\
Peru 2023 & 1 & 2013 & Advancing (ESG Integration) & 2023 \\
Philippines 2023 & 1 & 2013 & Advancing (ESG Integration) & 2023 \\
Poland 2019 & 0 & --- & None & 2019 \\
Portugal 2019 & 0 & --- & None & 2019 \\
Romania 2019 & 0 & --- & None & 2019 \\
Russia 2019 & 0 & --- & None & 2019 \\
Serbia 2019 & 0 & --- & None & 2019 \\
Slovak Republic 2019 & 0 & --- & None & 2019 \\
Slovenia 2019 & 0 & --- & None & 2019 \\
Tajikistan 2019 & 0 & --- & None & 2019 \\
Timor-Leste 2021 & 0 & --- & None & 2021 \\
Tunisia 2020 & 1 & 2019 & Unverified & 2020 \\
Türkiye 2019 & 1 & 2015 & Developing (Implementation) & 2019 \\
Ukraine 2019 & 0 & --- & None & 2019 \\
Uzbekistan 2019 & 0 & --- & None & 2019 \\
West Bank and Gaza 2019 & 0 & --- & None & 2019 \\
\bottomrule
\end{longtable}
}

\section{Bayesian model convergence diagnostics}
\label{app:diagnostics}

Convergence diagnostics for every focal parameter of the hierarchical model reported in the main
text are included directly in that table ($\hat{R}$ column); the full posterior summary also
includes $\text{ess}_{\text{bulk}}$ ranging from 797 (SBFN member) to 2{,}270 (credit $\times$
regulatory quality), comfortably above the conventional minimum threshold of 400 effective samples
per parameter for stable posterior summaries at four chains. No divergent transitions were flagged
during sampling (NUTS via the nutpie sampler, four chains of 1{,}000 post-warmup draws each after
1{,}000 tuning iterations).

\section{Global-sample SBFN status, regional summary}
\label{app:globalsbfn}

Of the 162 global-sample economy-years, 61 are coded as SBFN members as of their survey year and 101
as non-members. Table~\ref{tab:globalsbfnregion} reports the regional distribution of the 61
members, drawn from the SBFN data portal's regional classification; five of the 61 (Peru,
Philippines, Timor-Leste, Bangladesh, and Indonesia) carry the deeper per-country sourcing already
reported in Table~\ref{tab:appendixa} above, having been determined during the primary sample's
original research pass rather than from the portal roster alone.

\begin{table}[htbp]
\centering
\caption{Global-sample SBFN members, by SBFN region.}
\label{tab:globalsbfnregion}
\begin{tabular}{lc}
\toprule
Region & Member economy-years \\
\midrule
Latin America and Caribbean & 17 \\
Africa & 12 \\
Europe and Central Asia & 11 \\
East Asia and Pacific & 11 \\
South Asia & 6 \\
Middle East and North Africa & 4 \\
\midrule
Total & 61 \\
\bottomrule
\end{tabular}
\end{table}

\end{document}
